\documentclass[11pt]{article}

\usepackage[margin=1in]{geometry}
\usepackage{amsmath, amssymb}
\usepackage{xcolor}
\usepackage{listings}
\usepackage{hyperref}
\usepackage{array}
\usepackage{booktabs}

\hypersetup{
    colorlinks=true,
    linkcolor=blue!60!black,
    urlcolor=blue!60!black,
    citecolor=blue!60!black
}

\lstdefinestyle{bugpython}{
  language=Python,
  basicstyle=\ttfamily\small,
  keywordstyle=\color{blue!60!black},
  stringstyle=\color{green!40!black},
  commentstyle=\color{gray!70!black},
  showstringspaces=false,
  breaklines=true,
  frame=single,
  rulecolor=\color{gray!30},
  columns=fullflexible,
  keepspaces=true
}

\title{SymPy Bug Report}
\author{}
\date{}

\begin{document}

\maketitle

\section{Bug 21: solveset returns 1 for unsatisfiable principal-log equation}

\subsection{Minimal Reproducer}

\medskip

\begin{lstlisting}[style=bugpython]
import sympy
from sympy import Eq, I, S, log, pi, simplify, solveset, symbols

x = symbols("x")
eq = Eq(log(x), 2*pi*I)
sol = solveset(eq, x, domain=S.Complexes)

print("SymPy version:", sympy.__version__)
print("solveset result:", sol)
print("residual at returned value 1:",
      simplify(eq.lhs.subs(x, 1) - eq.rhs.subs(x, 1)))
\end{lstlisting}

\subsection{Observed Behavior}

\medskip

\begin{lstlisting}[style=bugpython]
SymPy version: 1.14.0
solveset result: {1}
residual at returned value 1: -2*I*pi
\end{lstlisting}

SymPy returns the finite solution set \(\{1\}\).  However, direct substitution
of the returned value into the original equation gives
\[
  \log(1) - 2\pi i = -2\pi i \ne 0.
\]
Thus the returned element does not satisfy the equation that \texttt{solveset}
was asked to solve.

\subsection{Expected Behavior}

The mathematically correct result is \(\varnothing\), represented in SymPy as
\texttt{EmptySet}.  There is no complex value \(x\) for which SymPy's principal
logarithm satisfies \(\log(x)=2\pi i\).

\subsection{Mathematical Explanation}

SymPy's complex \(\log\) is the principal logarithm.  For \(x\ne 0\), it can be
written as
\[
  \log(x) = \ln|x| + i\operatorname{Arg}(x),
  \qquad \operatorname{Arg}(x)\in(-\pi,\pi].
\]
Therefore every value in the image of the principal logarithm has imaginary
part in \((-\pi,\pi]\).  The target \(2\pi i\) has imaginary part \(2\pi\), so
it lies outside that image.

Exponentiating both sides gives \(x=e^{2\pi i}=1\), but this implication is not
reversible for the principal logarithm.  The value \(x=1\) satisfies
\(e^{\log(1)}=1\), but it does not satisfy the original equation because
\(\log(1)=0\), not \(2\pi i\).  The wrong answer is therefore not a harmless
branch representation or an equivalent form of a solution; it is a concrete
non-solution.

\subsection{Source-Level Diagnosis}

The diagnosis identifies the relevant code path in
\nolinkurl{sympy/solvers/solveset.py}, in the complex inversion helper
\texttt{\_invert\_complex} near line 550.  The traced call path is that
\texttt{solveset} rewrites \texttt{Eq(log(x), 2*pi*I)} as
\texttt{log(x) - 2*pi*I}, calls \texttt{\_solveset(..., \_check=True)}, and then
reaches the generic inversion branch.  That branch calls \texttt{\_invert} over
\texttt{S.Complexes}, which dispatches to \texttt{\_invert\_complex}.

For this equation, \texttt{\_invert\_complex(log(x) - 2*pi*I, \{0\}, x)}
first removes the additive constant and then handles \(\log(x)\) through the
generic function-inverse path.  It returns the transformed solution
\((x,\{e^{2\pi i}\})\), which simplifies to \((x,\{1\})\).  The problematic
source rule is the generic use of a function's inverse when one is available:

\begin{lstlisting}[style=bugpython]
if hasattr(f, 'inverse') and f.inverse() is not None and \
   not isinstance(f, TrigonometricFunction) and \
   not isinstance(f, HyperbolicFunction) and \
   not isinstance(f, exp):
    ...
    return _invert_complex(f.args[0],
                           imageset(Lambda(n, f.inverse()(n)), g_ys), symbol)
\end{lstlisting}

This treats \(\exp\) as a globally valid inverse for \(\log\) over the complex
numbers.  The missing condition is that the right-hand side must lie in the
principal branch range of \(\log\).  The later \texttt{\_check=True} filtering
identified in the diagnosis only performs a domain or singularity check through
\texttt{domain\_check}; it does not perform a residual check against the
original equation for the fully inverted finite set.  Consequently, the
candidate \(1\) is accepted even though substituting it into the original
equation gives a nonzero residual.  A plausible fix direction is to make the
complex inverse path for \(\log\) branch-range aware, or to return a conditional
set when the branch condition cannot be resolved.  The diagnosis also supports
a defensive residual check for finite candidates produced by branch-sensitive
inversion.

\subsection{Additional Failing Instances}

The same failure occurs for nonzero integer multiples of \(2\pi i\).  For
every nonzero integer \(n\) in the verified family, \(\operatorname{solveset}\)
is asked to solve
\[
  \log(x)=2\pi i n
\]
over the complex numbers.  SymPy returns \(\{1\}\), because
\(e^{2\pi i n}=1\), but \(2\pi n\) is outside the principal argument interval
\((-\pi,\pi]\).  Hence the expected solution set is \(\varnothing\) in each
case.

\begin{center}
\small
\renewcommand{\arraystretch}{1.3}
\begin{tabular}{@{}>{\raggedright\arraybackslash}p{0.44\linewidth}>{\raggedright\arraybackslash}p{0.23\linewidth}>{\raggedright\arraybackslash}p{0.23\linewidth}@{}}
\toprule
\textbf{Input / Expression} & \textbf{SymPy behavior} & \textbf{Expected behavior} \\
\midrule
\(\operatorname{solveset}(\log(x)=2\pi i,\ x,\ \mathbb{C})\) & returns \(\{1\}\); residual \(-2\pi i\) & \(\varnothing\) \\
\(\operatorname{solveset}(\log(x)=4\pi i,\ x,\ \mathbb{C})\) & returns \(\{1\}\); residual \(-4\pi i\) & \(\varnothing\) \\
\(\operatorname{solveset}(\log(x)=-2\pi i,\ x,\ \mathbb{C})\) & returns \(\{1\}\); residual \(2\pi i\) & \(\varnothing\) \\
\(\operatorname{solveset}(\log(x)=6\pi i,\ x,\ \mathbb{C})\) & returns \(\{1\}\); residual \(-6\pi i\) & \(\varnothing\) \\
\(\operatorname{solveset}(\log(x)=8\pi i,\ x,\ \mathbb{C})\) & returns \(\{1\}\); residual \(-8\pi i\) & \(\varnothing\) \\
\(\operatorname{solveset}(\log(x)=10\pi i,\ x,\ \mathbb{C})\) & returns \(\{1\}\); residual \(-10\pi i\) & \(\varnothing\) \\
\bottomrule
\end{tabular}
\end{center}

\subsection{Related Issues Assessment}

The bug-finding harness performed a best-effort deduplication check.  The
closest related public material found was documentation confirming that
SymPy's \(\log\) is principal-valued, documentation describing the contract of
\texttt{solveset}, and documentation for complex inversion in \texttt{solveset}.
No public issue, pull request, discussion, or release note was identified that
reports this exact reproducer or the exact wrong output \(\{1\}\).  The
deduplication verdict is therefore a new instance of a known failure mode: the
general need for principal-branch-aware inverse solving is known, but this
specific reproducible case is previously unreported and remains individually
actionable as a regression test and issue.

\subsection{Proposed SymPy Regression Test}

\medskip

\begin{lstlisting}[style=bugpython]
import pytest

from sympy import Eq, FiniteSet, I, S, log, pi, simplify, solveset, symbols


@pytest.mark.parametrize("n", [1, 2, -1, 3, 4, 5])
def test_solveset_principal_log_rejects_out_of_range_targets(n):
    x = symbols("x")
    eq = Eq(log(x), 2*pi*I*n)
    sol = solveset(eq, x, domain=S.Complexes)

    assert sol != FiniteSet(1)
    assert simplify(eq.lhs.subs(x, 1) - eq.rhs.subs(x, 1)) != 0
\end{lstlisting}

\end{document}
